\SourceLeaf{205}\begin{center}
\begin{tabular}{r|p{.82\linewidth}}
I & \\
 & Reversed\\
II & \\
 & Turned screw 7\\
III & \\
 & Turned screw 7 halfway back\\
IV & \\
 & Turned screw 3 if the spirit level is also to be adjusted.\\
V &
\end{tabular}
\end{center}

The deductions in reverse order are therefore as follows.

(1) $v^4=0$, because (2) $v^3=-v^2$ and equal small movements at 7 change $v$ by equal amounts.

(2) $v^3=-v^2$, because $v^2=u^1$, $v^3=-u^3$ and (3) $u^1=u^3$.

(3) $u^1=u^3$, because $S^1=S^3=u$ depends on $u$, hence $u^1=u^3$.

Since $S^1$ enters into the argument here, it is clear that before I the plane $T$ must already have been brought into a position in which the ends of the bubble are visible.

If the spirit level is now also to be adjusted, it must be shown that $i^5=0$.

This is apparent from the normal position of $S$, although that position evidently has no direct connection with $i=0$ and is connected only with $g=0$.

The synthesis is best begun from the original values.
\[
S = M+h+i+u,
\]
\[
\Sigma = M-h+i+u,
\]
\[
\begin{array}{rccc}
 & u & v & i \\
\cline{2-4}
{[\mathrm{I}]} & v^1 & v^1 & i^1 \\
{[\mathrm{II}]} & -v^1 & v^1 & i^1 \\
{[\mathrm{III}]} & v^1 & -v^1 & i^1 \\
{[\mathrm{IV}]} & 0 & 0 & i^1 \\
{[\mathrm{V}]} & 0 & 0 & 0 \\
\end{array}
\]

Among the 5 different positions for $v$, there is only 1\,case with $=+u$ and there are 4\,cases with $=-u$. The symbol $v$ can therefore be ignored entirely, and we may speak [only] of $u$.

